\documentclass[9pt,conference]{IEEEtran}
\usepackage{amssymb,amsthm,amsmath,array}
\usepackage{graphicx}
\usepackage[caption=false,font=footnotesize]{subfig}
\usepackage{xspace}
\usepackage[sort&compress, numbers]{natbib}
\usepackage{stmaryrd}
\usepackage{xcolor}
\usepackage{mathtools}
\usepackage{float}
\usepackage{textcomp}

\begin{document}
\title{\LaTeX~Author Guidelines for \\International Symposium on Computational Sensing}
\author{\IEEEauthorblockN{
        First Author\IEEEauthorrefmark{1}, 
        Second Author\IEEEauthorrefmark{2}, and 
        Third Author\IEEEauthorrefmark{3}
    }
    \IEEEauthorblockA{
        \IEEEauthorrefmark{1} First Author Affiliation.\\
        \IEEEauthorrefmark{2} Second Author Affiliation.\\
        \IEEEauthorrefmark{3} Third Author Affiliation.}
}
\maketitle
\begin{abstract}
    Abstract should contain approximately 100 to 150 words.
\end{abstract}

\section{Introduction}
Please follow the outlined formatting when submitting
your manuscript to the International Symposium on Computational Sensing (ISCS). This section provides guidelines for abstract submissions for all types of presentations in ISCS, \textit{i.e.,} platform, poster, and show-and-tell presentations. Specific requirements for show-and-tell abstracts are highlighted in Sec.~\ref{sec:show-and-tell}.
\subsection{Paper length}
The length of the paper must be up to 2 pages for technical content, figures, and references, and one optional third page containing only references.
\subsection{Font size}
Use a font size that is no smaller than 9 points throughout the paper, including figure captions.
\subsection{Bibliography}
List and number all bibliographical references at the end of
the paper. The references can be numbered in alphabetic order
or in order of appearance in the document.
When referring to
them in the text, type the corresponding reference number in
square brackets as shown at the end of this sentence \cite{code1}.

\section{Specific Guidelines for Show-and-Tell Abstract} \label{sec:show-and-tell}
The extended abstract for show-and-tell demos must follow the aforementioned format. The manuscript should 
\begin{itemize}
\item include pictures of the set-up;
\item include a clear explanation of the targeted application and computational method developed;
\item include a representative example of acquisitions using the set-up and its associated processing. 
\end{itemize}

\section{Conclusion}


\bibliographystyle{IEEEtran}
\bibliography{references}
\end{document}
